\title{Controlling Text-to-Image Diffusion by Orthogonal Finetuning}

\begin{abstract}
Large-scale text-to-image diffusion models demonstrate remarkable ability to produce photorealistic images from textual descriptions. An important unresolved question is how to effectively steer or control these powerful models for various downstream applications. To address this issue, we propose a principled finetuning approach—Orthogonal Finetuning~(OFT)—for adapting text-to-image diffusion models to specific downstream tasks. Unlike prior techniques, OFT can theoretically guarantee the preservation of hyperspherical energy, which captures the pairwise relationships among neurons on the unit hypersphere. We observe that maintaining this property is essential for retaining the semantic generation capabilities of text-to-image diffusion models. To further enhance finetuning stability, we introduce Constrained Orthogonal Finetuning (COFT), which applies an extra radius constraint on the hypersphere. We focus on two key downstream text-to-image finetuning tasks: subject-driven generation, where the objective is to produce subject-specific images from a few example images and a text prompt, and controllable generation, where the aim is to allow the model to incorporate supplementary control signals. Through empirical evaluation, we demonstrate that our OFT framework surpasses existing approaches in terms of both generation quality and convergence speed.
\end{abstract}

\section{Introduction}

Recent text-to-image diffusion models~\cite{saharia2022photorealistic,ramesh2022hierarchical,rombach2022high} have shown exceptional performance in producing high-fidelity images guided by text inputs. However, text prompts alone can remain ambiguous and may not provide sufficiently detailed or precise control over the resulting images. To overcome this limitation, this paper focuses on two categories of text-to-image generation tasks:

\begin{itemize}[leftmargin=*,nosep]

    \item \textbf{Subject-driven generation}~\cite{ruiz2023dreambooth}: Given only a handful of images depicting a subject, the task is to generate images featuring the same subject in varied contexts, guided by a text prompt.
    \item \textbf{Controllable generation}~\cite{zhang2023adding,mou2023t2i}: Provided with an extra control input (such as canny edges or segmentation maps), the goal is to create images consistent with both this control signal and a text prompt.
\end{itemize}

Both tasks fundamentally reduce to the challenge of how to efficiently finetune text-to-image diffusion models without compromising their generative capabilities developed during pretraining. We outline the criteria for an effective finetuning approach as follows: (1) \emph{training efficiency}: minimizing the number of trainable parameters and training epochs, and (2) \emph{generalizability preservation}: maintaining the model’s high-fidelity and diverse generative output. To achieve this, finetuning is commonly performed either by adjusting the neuron weights with a small learning rate (\emph{e.g.}, \cite{ruiz2023dreambooth}) or by introducing a minor component composed of re-parameterized neuron weights (\emph{e.g.}, \cite{hulora2022,zhang2023adding}). Despite their straightforwardness, neither of these finetuning approaches guarantees retention of the generative performance acquired during pretraining. Additionally, there is a lack of a systematic framework for designing effective finetuning strategies or for selecting appropriate hyperparameters, such as the optimal number of training epochs. A central challenge lies in the absence of a reliable metric to quantify how well the pretrained generative ability is preserved. Current finetuning methods implicitly assume that a smaller Euclidean distance between the finetuned and pretrained models corresponds to better preservation of pretrained capabilities. For this reason, finetuning typically proceeds with very small learning rates. Although this assumption may sometimes be valid, we contend that relying solely on the Euclidean difference from the pretrained model does not fully capture the extent of semantic retention. Consequently, employing a more structural metric to describe the disparity between the finetuned and pretrained models can substantially enhance both the preservation of pretrained performance and the stability of the finetuning process.

Motivated by empirical findings that hyperspherical similarity effectively captures semantic information~\cite{liu2017deep,liu2018decoupled,chen2020angular}, we employ hyperspherical energy~\cite{liu2018learning} to describe the pairwise relational structure among neurons. Hyperspherical energy is defined as the total sum of hyperspherical similarities (such as cosine similarity) between all pairs of neurons within the same layer, reflecting the degree of neuron uniformity on the unit hypersphere~\cite{liu2021learning}. We propose that an ideal finetuned model should exhibit minimal deviation in hyperspherical energy relative to the pretrained model. A straightforward approach would be to incorporate a regularizer to maintain constant hyperspherical energy during finetuning; however, there is no assurance that the difference in hyperspherical energy can be effectively minimized. To address this, we leverage an invariance property of hyperspherical energy — specifically, that the pairwise hyperspherical similarity remains exactly preserved when applying an identical orthogonal transformation to all neurons. Inspired by this invariance, we introduce Orthogonal Finetuning (OFT), which adapts large text-to-image diffusion models to downstream tasks without altering their hyperspherical energy. The key concept is to learn a shared orthogonal transformation at the layer level, ensuring the preservation of pairwise angles among neurons. OFT can also be interpreted as modifying the canonical coordinate system for neurons within the same layer. Furthermore, recognizing that a smaller Euclidean distance between the finetuned and pretrained models correlates with better retention of pretrained performance, we propose a variant called Constrained Orthogonal Finetuning (COFT), which restricts the finetuned model to lie within a hypersphere of fixed radius centered around the pretrained neurons.

The reasoning behind why orthogonal transformation is effective for finetuning neurons partially stems from the 2D Fourier transform, which allows an image to be broken down into magnitude and phase spectra. The phase spectrum, representing the angular relationship between the input and basis, retains most of the semantic content. For instance, an image’s phase spectrum combined with a random magnitude spectrum can still recreate the original image without losing its semantic meaning. This observation indicates that altering neuron orientations is crucial for semantically modifying the generated image, which aligns with the objectives of both subject-driven and controllable generation. However, allowing unrestricted changes in neuron directions can inevitably degrade the generative performance learned during pretraining. To limit this freedom, we propose maintaining the angle between any pair of neurons, based on the hypothesis that neuron angles are vital for encoding the knowledge of neural networks. Building on this insight, it is natural to learn a layer-shared orthogonal transformation for neurons within each layer to ensure that the hyperspherical energy remains constant.

We also take inspiration from orthogonal over-parameterized training~\cite{liu2021orthogonal}, which trains classification neural networks from scratch by orthogonally transforming a randomly initialized network. This approach leverages the fact that a randomly initialized neural network is expected to have a small hyperspherical energy, and \cite{liu2021orthogonal} aims to maintain this low hyperspherical energy during training (since lower energy is linked to better generalization in classification tasks~\cite{liu2018learning,lin2020regularizing}). Their work demonstrates that orthogonal transformations are sufficiently expressive to train neural networks that generalize well for classification problems. In contrast, our focus is on finetuning text-to-image diffusion models to enhance controllability and improve downstream generative performance. We highlight two main differences between OFT and \cite{liu2021orthogonal}. First, whereas \cite{liu2021orthogonal} is designed to minimize hyperspherical energy, OFT seeks to preserve the hyperspherical energy from the pretrained model to avoid disrupting its inherent structure during finetuning. Minimizing hyperspherical energy when finetuning diffusion models could compromise the original semantic structures. Second, OFT aims to minimize deviation from the pretrained model, resulting in a constrained variant, whereas \cite{liu2021orthogonal} does not impose such constraints. Finding a suitable balance between flexibility and stability is key to successful finetuning, and we argue that our OFT framework effectively accomplishes this balance. Our contributions are summarized as follows:

\begin{itemize}[leftmargin=*,nosep]

    \item We introduce a novel finetuning approach called Orthogonal Finetuning, designed to enhance the controllability of text-to-image diffusion models. To increase stability, we also present a constrained version that restricts the angular deviation from the pretrained model.
    \item Unlike existing finetuning techniques, OFT modifies the model while provably maintaining the hyperspherical energy, which we empirically identify as a crucial metric for preserving the generative semantics of the pretrained model.
    \item We apply OFT to two applications: subject-driven generation and controllable generation. Our extensive empirical evaluation shows notable improvements over previous methods in generation quality, convergence speed, and finetuning stability. Additionally, OFT demonstrates superior sample efficiency, achieving convergence with significantly fewer training images and epochs.
    \item For controllable generation, we propose a novel control consistency metric to assess controllability. The core concept is to estimate the control signal from the generated image and compare it to the original control input. This metric further confirms the strong controllability of OFT.
\end{itemize}

\section{Related Work}

\textbf{Text-to-image diffusion models}. Significant advancements~\cite{saharia2022photorealistic,ramesh2022hierarchical,rombach2022high,gu2022vector,nichol2022glide} have been achieved in text-to-image generation, largely driven by progress in diffusion-based generative models~\cite{ho2020denoising,song2021score,song2021denoising,dhariwal2021diffusion} and vision-language representation learning~\cite{su2020vlbert,lu2019vilbert,radford2021learning,wang2021simvlm,li2022blip,li2023blip,alayrac2022flamingo,schuhmann2022laion}. GLIDE~\cite{nichol2022glide} and Imagen~\cite{saharia2022photorealistic} train diffusion models in pixel space; GLIDE trains its text encoder jointly with a diffusion prior using paired text-image data, whereas Imagen uses a fixed pretrained text encoder. Stable Diffusion~\cite{rombach2022high} and DALL-E2~\cite{ramesh2022hierarchical} operate in latent space. Stable Diffusion employs VQ-GAN~\cite{esser2021taming} to learn a visual codebook serving as the latent space, while DALL-E2 leverages CLIP~\cite{radford2021learning} to create a joint latent embedding for images and text. Besides diffusion models, generative adversarial networks~\cite{reed2016generative,zhang2017stackgan,xu2018attngan,li2019controllable} and autoregressive models~\cite{ramesh2021zero,ding2021cogview,wu2022nuwa,yuscaling} have also been explored for text-to-image generation. OFT is fundamentally a model-agnostic finetuning strategy applicable to any text-to-image diffusion model.

\textbf{Subject-driven generation}. To avoid altering the subject, \cite{avrahami2022blended,nichol2022glide} incorporate user-provided masks as additional conditions. Inversion techniques~\cite{choi2021ilvr,dhariwal2021diffusion,ramesh2022hierarchical,gal2022image} enable context modification without changing the subject. \cite{hertz2022prompt} supports both local and global editing without requiring input masks. However, these methods often fail to maintain identity-specific subject details. Pivotal Tuning~\cite{roich2022pivotal} finetunes a generator around an initial inverted latent code with added regularization to preserve identity. Similarly, \cite{nitzan2022mystyle} learns a personalized generative face prior from multiple images of a person. \cite{casanova2021instance} can create diverse variations of an instance but may lose instance-specific features. DreamBooth~\cite{ruiz2023dreambooth} finetunes text-to-image diffusion models using a custom token and a limited set of subject images, employing a reconstruction loss along with a class-specific prior preservation loss. OFT builds on the DreamBooth framework, but rather than using standard finetuning with a small learning rate, it applies orthogonal transformations during finetuning.

\textbf{Controllable generation}. Image-to-image translation can be considered a type of controllable generation, and earlier approaches primarily rely on conditional generative adversarial networks~\cite{isola2017image,zhu2017toward,wang2018high,park2019semantic,choi2018stargan}. Diffusion models have also been employed for image-to-image translation tasks~\cite{saharia2022palette,voynov2022sketch,wang2022pretraining}. Recently, ControlNet~\cite{zhang2023adding} introduces a method to steer a pretrained diffusion model by finetuning and adapting it to incorporate additional control signals, achieving notable success in controllable generation. Similarly, another concurrent work, T2I-Adapter~\cite{mou2023t2i}, finetunes pretrained diffusion models to enhance control over the generated images. Adopting the same task framework as in \cite{zhang2023adding,mou2023t2i}, we apply OFT for finetuning pretrained diffusion models, consistently achieving improved controllability with fewer training samples and a smaller number of finetuning parameters. Importantly, OFT does not add any extra computational cost during inference.

\textbf{Model finetuning}. The practice of finetuning large pretrained models on downstream tasks has become increasingly common~\cite{devlin2018bert,brown2020language,he2022masked}. Within this context, adaptation techniques (\emph{e.g.}, \cite{hulora2022,houlsby2019parameter,pfeiffer2020adapterfusion}) have been extensively explored in natural language processing. LoRA~\cite{hulora2022} is closely related to OFT and assumes a low-rank structure for the additive weight updates during finetuning. In contrast, OFT employs a layer-shared orthogonal transformation to modify neuron weights multiplicatively, which provably maintains the pairwise angles among neurons within the same layer, leading to enhanced stability.

\section{Orthogonal Finetuning}

\subsection{Why Does Orthogonal Transformation Make Sense?}

We begin by examining why orthogonal transformation is a sensible choice for finetuning text-to-image diffusion models. This inquiry is divided into two parts: (1) the reason for finetuning the angles of neurons (\emph{i.e.}, their directions), and (2) the rationale behind using orthogonal transformation to adjust these angles.

Regarding the first question, we take inspiration from the empirical findings in \cite{liu2018decoupled,chen2020angular} that angular feature differences effectively represent the semantic gap. SphereNet~\cite{liu2017deep} demonstrates that training a neural network with all neurons normalized onto a unit hypersphere achieves comparable capacity and even enhances generalizability, suggesting that neuron directions alone can capture the most critical information from the data. To further illustrate the significance of neuron angles, we conduct a simple experiment shown in Figure~\ref{fig:toy_ae}, where we train a standard convolutional autoencoder on a set of flower images. During training, we employ the conventional inner product to generate the feature map (where $z$ represents the element output of the convolution kernel $\bm{w}$ and $\bm{x}$ is the input within the sliding window). At test time, we compare three methods of producing the feature map: (a) the inner product used in training, (b) the magnitude information, and (c) the angular information. The results depicted in Figure~\ref{fig:toy_ae} demonstrate that the angular information of neurons nearly perfectly reconstructs the input images, whereas the neuron magnitudes carry no meaningful information. It is important to note that we do not apply the cosine activation (c) during training; training is solely based on the inner product. These findings suggest that neuron angles (directions) play a dominant role in encoding the semantic content of input images. Consequently, adjusting neuron directions is likely to be a more effective approach for altering semantic information.

Regarding the second question, the most straightforward method to fine-tune neuron directions is to rotate or reflect all neurons simultaneously within the same layer, naturally introducing orthogonal transformations. While it might be possible to apply more flexible angular transformations that rotate different neurons by distinct angles, we find orthogonal transformations offer an ideal balance between flexibility and structure. Additionally, \cite{liu2021orthogonal} shows that orthogonal transformations are sufficiently powerful for neural network learning. To support this claim, we conduct an experiment demonstrating the effective regularization imposed by the orthogonality constraint. We carry out a controllable generation experiment under the ControlNet~\cite{zhang2023adding} setup, with results shown in Figure~\ref{fig:ortho}. We observe that our standard OFT method performs robustly and achieves precise control after training completes (epoch 20). In contrast, OFT without the orthogonality constraint fails to generate any realistic images and exhibits no control capability. This experiment confirms the critical role of the orthogonality constraint in OFT.

\subsection{General Framework}

The traditional finetuning approach generally employs gradient descent with a low learning rate to update a model (or selected layers within a model). This small learning rate implicitly restricts the changes from the pretrained model, so standard finetuning essentially tries to train the model while implicitly minimizing $\thickmuskip=2mu \medmuskip=2mu \| \bm{M} - \bm{M}^0\|$, where $\bm{M}$ represents the finetuned model weights and $\bm{M}^0$ denotes the pretrained model weights. While this implicit limitation is intuitive, it may still allow excessive flexibility when finetuning a large model. To counter this, LoRA introduces an extra low-rank restriction on the weight updates, i.e., $\thickmuskip=2mu \medmuskip=2mu \text{rank}(\bm{M} - \bm{M}^0)=r'$, where $r'$ is set to a small value. Unlike LoRA, OFT imposes a constraint on the pairwise neuron similarity: $\thickmuskip=2mu \medmuskip=2mu \| \text{HE}(\bm{M})-\text{HE}(\bm{M}^0) \|=0$, where $\text{HE}(\cdot)$ stands for the hyperspherical energy of a weight matrix. 

As an example, consider a fully connected layer $\thickmuskip=2mu \medmuskip=2mu \bm{W}=\{\bm{w}_1,\cdots,\bm{w}_n\}\in\mathbb{R}^{d\times n}$, where each $\bm{w}_i\in\mathbb{R}^d$ is the $i$-th neuron (with pretrained weights $\bm{W}^0$). The output vector $\thickmuskip=2mu \medmuskip=2mu \bm{z}\in\mathbb{R}^n$ of this layer is computed as $\thickmuskip=2mu \medmuskip=2mu \bm{z}=\bm{W}^\top\bm{x}$, with $\thickmuskip=2mu \medmuskip=2mu \bm{x}\in\mathbb{R}^d$ being the input vector. OFT can be viewed as minimizing the difference in hyperspherical energy between the finetuned model and the pretrained model:

\begin{equation}\label{eq:oft_principle}
\footnotesize
\min_{\bm{W}} \left\| \text{HE}(\bm{W})-\text{HE}(\bm{W}^0)\right\|~~\Leftrightarrow~~\min_{\bm{W}} \bigg{\|} \sum_{i\neq j} \|\hat{\bm{w}}_i-\hat{\bm{w}}_j\|^{-1}-\sum_{i\neq j} \|\hat{\bm{w}}^0_i-\hat{\bm{w}}^0_j\|^{-1}\bigg{\|}
\end{equation}

Here, $\thickmuskip=2mu \medmuskip=2mu \hat{\bm{w}}_i := \bm{w}_i/\|\bm{w}_i\|$ is the normalized $i$-th neuron, and the hyperspherical energy of layer $\bm{W}$ is defined as $\thickmuskip=2mu \medmuskip=2mu \text{HE}(\bm{W}) := \sum_{i\neq j} \|\hat{\bm{w}}_i - \hat{\bm{w}}_j\|^{-1}$. It is straightforward to see that the minimal achievable value of Eq.~\eqref{eq:oft_principle} is zero. This minimum is attainable if $\bm{W}$ and $\bm{W}^0$ differ only by a rotation or reflection, i.e., $\thickmuskip=2mu \medmuskip=2mu \bm{W} = \bm{R} \bm{W}^0$, where $\thickmuskip=2mu \medmuskip=2mu \bm{R} \in \mathbb{R}^{d \times d}$ is an orthogonal matrix (with determinant $1$ or $-1$ corresponding to rotation or reflection, respectively). This encapsulates the essence of OFT: finetuning the neural network by learning layer-shared orthogonal matrices that transform neurons in

\begin{equation}\label{eq:oft_general}
    \footnotesize
    \bm{z}=\bm{W}^\top\bm{x}=(\bm{R}\cdot\bm{W}^0)^\top\bm{x},~~~\text{s.t.}~\bm{R}^\top\bm{R}=\bm{R}\bm{R}^\top=\bm{I}
\end{equation}

Here, $\bm{W}$ represents the OFT-finetuned weight matrix and $\bm{I}$ stands for the identity matrix. OFT is depicted in Figure~\ref{fig:block}. Similar to zero initialization in LoRA, it is crucial to ensure that OFT fine-tunes the pretrained model precisely from $\bm{W}^0$. To accomplish this, we initialize the orthogonal matrix $\bm{R}$ as an identity matrix so that the fine-tuned model begins with the pretrained weights. To maintain the orthogonality of matrix $\bm{R}$, differential orthogonalization techniques as described in \cite{liu2021orthogonal,lezcano2019cheap} can be employed. We will further elaborate on how to preserve orthogonality efficiently in computation.

\subsection{Efficient Orthogonal Parameterization}\label{sect:eff_ortho}

Although standard orthogonalization methods like the Gram-Schmidt process are differentiable, they tend to be computationally expensive in practice~\cite{liu2021orthogonal}. To improve efficiency, we utilize the Cayley parameterization to generate the orthogonal matrix. In particular, the orthogonal matrix is constructed as $\thickmuskip=2mu \medmuskip=2mu \bm{R}=(\bm{I}+\bm{Q})(\bm{I}-\bm{Q})^{-1}$, where $\bm{Q}$ is a skew-symmetric matrix satisfying $\thickmuskip=2mu \medmuskip=2mu \bm{Q}=-\bm{Q}^\top$. This approach provides greater efficiency at the cost of a slight limitation—Cayley parameterization produces orthogonal matrices with determinant $1$, thus belonging to the special orthogonal group. Fortunately, this constraint does not impact practical performance. Even when using the Cayley transform to parameterize $\bm{R}$, it can remain parameter-inefficient for large $d$. To tackle this issue, we propose representing $\bm{R}$ as a block-diagonal matrix composed of $r$ blocks, resulting in:

\begin{equation}
    \footnotesize
    \bm{R}=\text{diag}(\bm{R}_1,\bm{R}_2,\cdots,\bm{R}_r)=\begin{bmatrix}
    \bm{R}_1\in \text{O}(\frac{d}{r}) &  &\\
    &\ddots & \\
    & & \bm{R}_r\in \text{O}(\frac{d}{r})
    \end{bmatrix}\in \text{O}(d)
\end{equation}

where $\text{O}(d)$ represents the orthogonal group in $d$ dimensions, and $\thickmuskip=2mu \medmuskip=2mu \bm{R}\in\mathbb{R}^{d\times d}$ and $\thickmuskip=2mu \medmuskip=2mu \bm{R}_i\in\mathbb{R}^{d/r\times d/r}, \forall i$ denote orthogonal matrices. When $\thickmuskip=2mu \medmuskip=2mu r=1$, the block-diagonal orthogonal matrix reduces to a conventional unconstrained orthogonal matrix. For an orthogonal matrix of size $\thickmuskip=2mu \medmuskip=2mu d \times d$, the total number of parameters is $\thickmuskip=2mu \medmuskip=2mu d(d-1)/2$, which leads to a computational complexity of $\mathcal{O}(d^2)$. For an $r$-block diagonal orthogonal matrix, the parameter count is $\thickmuskip=2mu \medmuskip=2mu d(d/r - 1)/2$, yielding a complexity of $\mathcal{O}(d^2 / r)$. Optionally, we can share the block matrices across blocks, i.e., $\thickmuskip=2mu \medmuskip=2mu \bm{R}_i = \bm{R}_j$ for all $i \neq j$, which further lowers the number of parameters, reducing the complexity to $\mathcal{O}(d^2 / r^2)$. Despite these parameter-efficiency improvements, the resulting matrix $\bm{R}$ remains orthogonal, ensuring that there is no compromise in maintaining hyperspherical energy.

We analyze how OFT compares with LoRA regarding parameter efficiency. For LoRA with a low-rank parameter $r'$, the count of trainable parameters is $\thickmuskip=2mu \medmuskip=2mu r'(d + n)$. If both $r$ and $r'$ scale with the neuron dimension $d$ (e.g., $\thickmuskip=2mu \medmuskip=2mu r = r' = \alpha d$ for some constant $\alpha$ with $0 < \alpha \leq 1$), then the parameter complexity for LoRA is $\mathcal{O}(d^2 + dn)$, whereas for OFT it is $\mathcal{O}(d)$. To illustrate this difference concretely, consider a weight matrix sized $\thickmuskip=2mu \medmuskip=2mu 128 \times 128$: LoRA requires $2,048$ trainable parameters with $\thickmuskip=2mu \medmuskip=2mu r' = 8$, whereas OFT requires only $960$ trainable parameters with $\thickmuskip=2mu \medmuskip=2mu r=8$ (without block sharing).

\subsection{Constrained Orthogonal Finetuning}

We can further restrict the flexibility of the original OFT by confining the finetuned model to lie within a small vicinity of the pretrained model. In particular, COFT performs the forward pass as follows:

\begin{equation}\label{eq:coft_general}
    \footnotesize
    \bm{z} = \bm{W}^\top \bm{x} = (\bm{R} \cdot \bm{W}^0)^\top \bm{x}, \quad \text{s.t.} \quad \bm{R}^\top \bm{R} = \bm{R} \bm{R}^\top = \bm{I}, \quad \norm{\bm{R} - \bm{I}} \leq \epsilon
\end{equation}

which enforces both an orthogonality constraint and an $\epsilon$-radius constraint relative to the identity matrix. The orthogonality constraint can be imposed using the Cayley parameterization introduced in Section~\ref{sect:eff_ortho}. However, integrating the $\epsilon$-deviation constraint directly into the Cayley-parameterized orthogonal matrix is nontrivial. To gain further understanding of the Cayley transform, we apply the Neumann series to approximate $\thickmuskip=2mu \medmuskip=2mu \bm{R} = (\bm{I} + \bm{Q})(\bm{I} - \bm{Q})^{-1}$ as $\thickmuskip=2mu \medmuskip=

(the series converges in the operator norm). Consequently, we can incorporate the constraint $\thickmuskip=2mu \medmuskip=2mu\|\bm{R}-\bm{I}\|\leq\epsilon$ within the Cayley transform, leading to an equivalent constraint $\thickmuskip=2mu \medmuskip=2mu \|\bm{Q}-\bm{0}\|\leq\epsilon'$, where $\bm{0}$ represents an all-zero matrix and $\epsilon'$ is another error hyperparameter (distinct from $\epsilon$). This new constraint on the matrix $\bm{Q}$ can be conveniently imposed via projected gradient descent. To initialize the orthogonal matrix $\bm{R}$ as the identity, we set $\bm{Q}$ initially to an all-zero matrix. COFT can be interpreted as combining two explicit constraints: minimizing the hyperspherical energy difference and limiting deviation from the pretrained model. The latter constraint is commonly implicit in existing finetuning approaches, but COFT explicitly enforces it. Although OFT exhibits excellent performance, we find that COFT offers even greater finetuning stability due to this explicit deviation constraint. Figure~\ref{fig:eps_coft} illustrates how the choice of $\epsilon$ influences COFT's performance. It can be seen that $\epsilon$ regulates the finetuning flexibility. Larger $\epsilon$ values produce COFT-finetuned models closely resembling OFT-finetuned ones, whereas smaller $\epsilon$ values yield COFT-finetuned models behaving more like the pretrained text-to-image diffusion model.

\subsection{Re-scaled Orthogonal Finetuning}

We introduce a straightforward enhancement to the original OFT by also learning a scaling coefficient for each neuron’s magnitude. This approach is motivated by the observation that rescaling neurons does not affect hyperspherical energy, since magnitude normalization removes it. Specifically, we define the forward pass as: $\thickmuskip=2mu \medmuskip=2mu\bm{z}=(\bm{R}\bm{W}^0\bm{D})^\top\bm{x}$\footnote{\textbf{Errata}: In the NeurIPS camera ready version, the forward pass of re-scaled OFT was incorrectly written as $\thickmuskip=2mu \medmuskip=2mu\bm{z}=(\bm{D}\bm{R}\bm{W}^0)^\top\bm{x}$. The original implementation remains correct, so the results in Appendix~\ref{app:rescaled} are unaffected.} where $\thickmuskip=2mu \medmuskip=2mu \bm{D}=\text{diag}(s_1,\cdots,s_n)\in\mathbb{R}^{n\times n}$ is a learnable diagonal matrix with strictly positive diagonal entries $s_1,\cdots,s_n$. Unlike the original OFT forward pass in Eq.~\eqref{eq:oft_general}, where only $\bm{R}$ is learnable, here both the diagonal matrix $\bm{D}$ and the orthogonal matrix $\bm{R}$ are trainable. This re-scaled OFT variant increases OFT’s flexibility with a negligible increase in parameters. Our experiments focus on the original OFT to highlight the efficacy of orthogonal transformation alone, but we note that re-scaled OFT generally performs better (see Appendix~\ref{app:rescaled}).

\section{Intriguing Insights and Discussions}

\textbf{OFT is architecture-agnostic}. In principle, OFT can be applied to any neural network architecture. For Transformers, LoRA typically targets the attention weights~\cite{hulora2022}. To ensure fair comparisons with LoRA, we restrict OFT to finetune attention weights in our experiments. Beyond fully connected layers, OFT is also well-suited for finetuning convolutional layers, since the block-diagonal structure of $\bm{R}$ carries meaningful interpretations in convolutional layers (unlike LoRA). When the number of blocks equals the number of input channels, each block acts on a single neuron channel, resembling depth-wise convolution kernels~\cite{chollet2017xception}. When all blocks in $\bm{R}$ are shared, OFT applies a shared orthogonal transformation across channels. We provide a preliminary investigation on finetuning convolutional layers with OFT in Appendix~\ref{app:convolution}.

\textbf{Relation to LoRA}. LoRA adds a low-rank matrix update to the pretrained weight matrix, which helps to maintain the original information by avoiding significant changes. In contrast, OFT applies an orthogonal (full-rank) transformation to the pretrained weight matrix, ensuring that the pretraining knowledge is preserved. The forward pass of OFT can be expressed as $\thickmuskip=2mu \medmuskip=2mu \bm{z}=(\bm{R}\bm{W}^0)^\top\bm{x}=(\bm{W}^0+(\bm{R}-\bm{I})\bm{W}^0)^\top\bm{x}$, where $\thickmuskip=2mu \medmuskip=2mu (\bm{R}-\bm{I})\bm{W}^0$ resembles the low-rank weight update used by LoRA. Given that $\bm{W}^0$ is generally full-rank, OFT also effectively performs a low-rank weight update when $\thickmuskip=2mu \medmuskip=2mu \bm{R}-\bm{I}$ has low rank. Similar to LoRA, which employs a rank parameter $r'$, OFT uses a diagonal block parameter $r$ to limit the number of trainable parameters. Notably, LoRA and OFT represent two distinct strategies for parameter efficiency: LoRA leverages low-rank structure to reduce trainable parameters, whereas OFT exploits sparsity through block-diagonal orthogonality.

\textbf{Reasons for OFT’s faster convergence}. First, as illustrated in Figure~\ref{fig:toy_ae}, the most impactful update for altering semantics involves changing neuron directions, which is precisely the purpose OFT serves. Additionally, OFT can be interpreted as finetuning neurons constrained to a smooth hyperspherical manifold, leading to a more favorable optimization landscape. This advantage has also been empirically demonstrated in \cite{liu2021orthogonal}.

\textbf{Rationale for not minimizing hyperspherical energy}. Unlike \cite{liu2021orthogonal}, our approach does not seek to minimize hyperspherical energy. In classification tasks, neurons should avoid redundancy. Minimizing hyperspherical energy results in neurons being evenly distributed on the hypersphere, which is not a suitable goal for finetuning, as it risks erasing valuable pretraining information.

\textbf{Trade-off between flexibility and regularity in finetuning}. We identify a fundamental trade-off between flexibility and regularity. Standard finetuning offers the highest flexibility but suffers from poor stability and a tendency for model collapse. Surprisingly straightforward, OFT achieves a favorable balance between flexibility and regularity by maintaining the pairwise angles between neurons. Additionally, the block-diagonal parameterization can be interpreted as a more stringent regularization of the orthogonal matrix.

\textbf{No additional inference overhead}. In contrast to ControlNet, our OFT framework does not add any extra inference cost to the finetuned model. During inference, we simply multiply the learned orthogonal matrix $\bm{R}$ with the pretrained weight matrix $\bm{W}^0$ to obtain an equivalent weight matrix $\thickmuskip=2mu \medmuskip=2mu \bm{W}=\bm{R}\bm{W}^0$. Consequently, the inference speed remains identical to that of the pretrained model.

\section{Experiments and Results}

\textbf{General settings}. In our experiments, we utilize Stable Diffusion v1.5~\cite{rombach2022high} as the pretrained text-to-image model. To ensure fairness, we randomly select generated images from each method. For subject-driven generation, we largely follow DreamBooth~\cite{ruiz2023dreambooth}. For controllable generation, we follow the protocols of ControlNet~\cite{zhang2023adding} and T2I-Adapter~\cite{mou2023t2i}. To maintain fairness when comparing to LoRA, we apply OFT or COFT only to the same layers where LoRA is employed. Additional results and details can be found in Appendix~\ref{app:settings}.

\subsection{Subject-driven Generation}

\textbf{Settings}. We use DreamBooth~\cite{ruiz2023dreambooth} and LoRA~\cite{hulora2022} as baseline methods. All methods use the same loss function as DreamBooth. For DreamBooth and LoRA, we generally follow their original papers and employ their best hyperparameter configurations. Further results are included in Appendices~\ref{app:settings}, \ref{app:coft_oft}, \ref{app:more_qualitative}, and \ref{app:failure}.

\textbf{Stability and convergence of finetuning}. Initially, we assess the finetuning stability and convergence rate of DreamBooth, LoRA, OFT, and COFT. The findings are presented in Figure~\ref{fig:teaser} and Figure~\ref{fig:db_convergence}. It is evident that both COFT and OFT can finetune the diffusion model with considerable stability. After 400 iterations, DreamBooth and OFT variants demonstrate effective control, whereas LoRA does not maintain the subject’s identity. At 2000 iterations, DreamBooth begins producing degenerate images, and LoRA fails to generate a yellow shirt, instead producing yellow fur. Conversely, OFT and COFT continue to provide stable and consistent control over the generated outputs. These observations confirm the rapid yet stable convergence capability of our OFT framework in subject-driven generation. We emphasize that the robustness to the number of finetuning iterations is crucial, as it alleviates the difficulty of tuning iteration counts for different subjects. For both OFT and COFT, a relatively large iteration count can be set directly without meticulous adjustment. Specifically for COFT with an appropriate $\epsilon$, setting the learning rate and iteration number becomes straightforward and effortless.

\textbf{Quantitative evaluation}. Following \cite{ruiz2023dreambooth}, we perform a quantitative comparison to measure subject fidelity (DINO~\cite{caron2021emerging}, CLIP-I~\cite{radford2021learning}), text prompt fidelity (CLIP-T~\cite{radford2021learning}), and sample diversity (LPIPS~\cite{zhang2018unreasonable}). CLIP-I calculates the average pairwise cosine similarity between CLIP embeddings of generated and real images. DINO is analogous to CLIP-I, but employs ViT S/16 DINO embeddings instead. CLIP-T is the average cosine similarity between the CLIP embeddings of the text prompt and the generated images. We further assess the average LPIPS cosine similarity among generated images for the same subject under the identical text prompt. Table~\ref{table:dreambooth_final} demonstrates that both COFT and OFT substantially outperform DreamBooth and LoRA on the DINO and CLIP-I metrics, while delivering slightly better or comparable results in prompt fidelity and diversity measures. To reduce randomness, each method finetunes the same pretrained model 30 times with different random seeds. The outcomes indicate that our OFT framework not only provides superior convergence and stability but also consistently achieves improved final performance.

\textbf{Qualitative comparison}. To better grasp the advantages of OFT, we present several randomly selected examples for subject-driven generation in Figure~\ref{fig:dreambooth_final}. For a fair evaluation, all examples are generated from the identical finetuned model for each method, ensuring no text prompt is individually optimized for the final outputs. For each approach, we pick the model that attains the highest validation CLIP scores. As shown in Figure~\ref{fig:dreambooth_final}, both OFT and COFT demonstrate strong preservation of semantic subject features, whereas LoRA frequently fails to maintain the subject’s identity (e.g., LoRA completely loses the subject identity in the bowl example). Additionally, OFT and COFT provide much more precise control through text prompts, while DreamBooth, although it preserves the subject identity, often does not generate images that adhere well to the text prompts (e.g., the first row of the bowl example). This qualitative comparison illustrates that our OFT framework achieves superior controllability and subject preservation simultaneously. Furthermore, OFT’s performance is robust to the number of iterations, performing well even with a high iteration count, unlike DreamBooth or LoRA. Additional qualitative examples are provided in Appendix~\ref{app:more_qualitative}. We also conduct a human evaluation in Appendix~\ref{app:human}, which further confirms the advantages of OFT.

\subsection{Controllable Generation}

\textbf{Settings}. We utilize ControlNet~\cite{zhang2023adding}, T2I-Adapter~\cite{mou2023t2i}, and LoRA~\cite{hulora2022} as baseline methods. The main paper examines three demanding controllable generation tasks: Canny edge to image (C2I) on the COCO dataset~\cite{lin2014microsoft}, segmentation map to image (S2I) on the ADE20K dataset~\cite{zhou2017scene}, and landmark to face (L2F) on the CelebA-HQ dataset~\cite{karras2018progressive,xia2021tedigan}. Each method finetunes Stable Diffusion (SD) v1.5 on these datasets for 20 epochs. Further results are available in Appendix~\ref{app:more_qualitative}, \ref{app:more_control_task}, and \ref{app:failure}.

\textbf{Convergence}. We assess the convergence rates of ControlNet, T2I-Adapter, LoRA, and COFT on the L2F task, providing both quantitative and qualitative analyses. For the evaluation metric, we calculate the mean $\ell_2$ distance between the control face landmarks and the predicted face landmarks. In Figure~\ref{fig:face_convergence}, we illustrate the face landmark error for models fine-tuned over varying numbers of epochs. It is evident that both COFT and OFT demonstrate significantly faster convergence. LoRA requires 20 epochs to reach the performance level that our OFT framework attains by the 8th epoch. We highlight that OFT and COFT involve a comparable number of trainable parameters to LoRA (considerably fewer than ControlNet), yet converge much more efficiently than the existing methods. Furthermore, the rapid convergence of OFT is corroborated by the results shown in Figure~\ref{fig:teaser}. The right-hand example in Figure~\ref{fig:teaser} indicates that OFT is far more data-efficient than ControlNet and LoRA, as it achieves good convergence using only 5\% of the ADE20K dataset. For the qualitative comparison, we concentrate on OFT, COFT, and ControlNet, since ControlNet attains landmark errors closest to our methods. As presented in Figure~\ref{fig:face_convergence_qual}, both OFT and COFT converge steadily, with the generated face pose progressively aligning with the control landmarks. Unlike our smooth and stable convergence, ControlNet’s controllability abruptly appears after the 8th epoch, mirroring the sudden convergence behavior reported in \cite{zhang2023adding}. This stability in convergence renders the OFT framework more user-friendly in practical applications, as its training dynamics are considerably smoother and more predictable. Consequently, identifying optimal hyperparameters for OFT is easier.

\textbf{Quantitative comparison}. We propose a control consistency metric to assess the effectiveness of controllable generation. The core concept is to extract the control signal from the generated image and then compare it against the original input control signal. For the C2I task, we calculate IoU and F1 score. For the S2I task, we measure mean IoU, mean accuracy, and overall accuracy. For the L2F task, we determine the mean $\ell_2$ distance between the control landmarks and the predicted landmarks. Additional details about these consistency metrics can be found in Appendix~\ref{app:settings}. For all methods compared, we utilize the optimal hyperparameter configurations. The results presented in Table~\ref{table:control_gen} indicate that both OFT and COFT achieve significantly stronger and more precise control compared to other methods. We note that adapter-based methods (e.g., T2I-Adapter and ControlNet) exhibit slower convergence and produce inferior final outcomes. Relative to ControlNet, LoRA performs better on the S2I task but underperforms on the C2I and L2F tasks. Overall, we observe that LoRA’s performance ceiling remains relatively low despite careful hyperparameter tuning. In contrast, our OFT framework has not yet reached its performance limit, as we empirically observe continued improvements with an increasing number of trainable parameters. We highlight that our quantitative evaluation approach for controllable generation is a novel contribution, as it enables precise assessment of the control performance of fine-tuned models (subject to the accuracy of the utilized off-the-shelf segmentation/detection model).

\textbf{Qualitative comparison}. We further conduct a qualitative comparison of OFT and COFT with leading state-of-the-art approaches, such as ControlNet, T2I-Adapter, and LoRA. The randomly generated images displayed in Figure~\ref{fig:control_final} indicate that OFT and COFT not only produce high-quality, realistic images but also provide precise control. In the S2I task, it is evident that LoRA completely fails to generate images that conform to the input segmentation map, whereas ControlNet, OFT, and COFT successfully manage the generation process. Compared to ControlNet, both OFT and COFT generate higher-fidelity images featuring richer details and more accurate geometric structures, all while using significantly fewer model parameters. For the C2I task, both OFT and COFT can create semantically coherent images based on coarse Canny edges, whereas T2I-Adapter and LoRA exhibit much poorer performance. In the L2F task, our approach achieves the most precise pose control for generated faces, even in difficult facial poses. Across all three control tasks, we demonstrate that OFT and COFT yield qualitatively superior images compared to existing state-of-the-art methods, showcasing the efficacy of the OFT framework in controllable generation. To offer a more thorough qualitative evaluation, additional examples for all three control tasks are provided in Appendix~\ref{app:control_exp}. Furthermore, we show that OFT performs well on a broader range of control tasks, including dense pose to human body, sketch to image, and depth to image, as detailed in Appendix~\ref{app:more_control_task}.

\section{Concluding Remarks and Open Problems}

Inspired by the insight that angular relationships between neurons play a vital role in defining visual semantics, we introduce a straightforward yet powerful finetuning approach—orthogonal finetuning (OFT)—for controlling text-to-image diffusion models. Our work focuses on two specific applications: subject-driven generation and controllable generation. Compared with prior techniques, OFT achieves enhanced controllability and finetuning stability while requiring fewer finetuning parameters. Crucially, OFT does not incur any additional inference cost, enabling efficient deployment of the model.

OFT also raises several intriguing open questions. First, the orthogonality in OFT is ensured through Cayley parametrization, which necessitates computing a matrix inverse. This imposes a slight constraint on the scalability of OFT. Although we mitigate this through block diagonal parametrization, finding ways to accelerate the matrix inverse computation in a differentiable manner remains a challenge. Second, OFT offers distinctive potential for compositionality, as the orthogonal matrices derived from multiple OFT finetuning tasks can be multiplied while still yielding an orthogonal matrix. Investigating whether this collection of orthogonal matrices retains the knowledge from all downstream tasks is an engaging avenue for future research. Lastly, the parameter efficiency of OFT heavily relies on the block diagonal design, which inevitably introduces some bias and restricts flexibility. Developing methods to enhance parameter efficiency in a more effective and less biased manner remains an important open problem.

\end{document}